\documentclass[11pt]{article}
\usepackage[T1]{fontenc}
\usepackage{amsmath,amssymb,amsthm,mathtools}
\usepackage{graphicx}
\usepackage[margin=1in]{geometry}
\usepackage{microtype}
\usepackage[hidelinks]{hyperref}
\usepackage{xcolor}

\newtheorem{theorem}{Theorem}
\newtheorem{lemma}[theorem]{Lemma}
\newtheorem{remark}[theorem]{Remark}
\newcommand{\T}{\mathfrak T}
\newcommand{\Y}{\mathfrak Y}
\newcommand{\F}{\mathfrak F}
\newcommand{\B}{\mathfrak B}
\newcommand{\cH}{\mathcal H}
\newcommand{\id}{\operatorname{Id}}
\newcommand{\Tr}{\operatorname{Tr}}
\newcommand{\Ran}{\operatorname{Ran}}
\newcommand{\cb}{\mathrm{cb}}

\title{The ECD completion of Hermitian maps is generated by its CP cone\\
\large Candidate full proof of the conjecture in arXiv:1810.10922}
\author{Autonomous proof packet --- agent\_lane\_12}
\date{11 August 2026}

\begin{document}
\maketitle

\begin{abstract}
Let $G$ be a positive discrete unbounded operator.  Shirokov conjectured
that every Hermitian-preserving map in the energy-constrained-diamond (ECD)
completion $\Y_G(A,B)$ is a difference of two maps in the completed
completely-positive cone $\F_G^0(A,B)$.  We prove the conjecture.  Fixing
$E_0>0$ and writing $R=(I+G/E_0)^{-1/2}$, the weighted map
$\widehat\Phi(\rho)=\Phi(R\rho R)$ is an ordinary completely bounded map,
and its diamond norm is equivalent to the ECD norm, with constants $1$ and
$2$.  Discreteness of $G$ makes $R$ compact.  Hence the normal dual of
$\widehat\Phi$ is a cb-norm limit of normal Hermitian maps taking values in
finite spectral corners.  A summable telescoping approximation, normal
Wittstock decomposition in each corner, and conjugation by the bounded
finite-corner inverse of $R$ produce two ECD-convergent series of bounded CP
maps.  Their limits lie in $\F_G^0(A,B)$ and their difference is $\Phi$.
\end{abstract}

\section{Source problem}

The source is M.~E.~Shirokov, \emph{On completion of the cone of CP linear
maps with respect to the energy-constrained diamond norm},
arXiv:1810.10922v2 (2019)~\cite{Shirokov}.  On page 23 it asks:

\begin{center}
\includegraphics[width=0.98\textwidth]{figures/open_problem_crop.png}
\end{center}

In the source notation, $\T_G(\cH_A)$ is the linear span of positive
trace-class operators of finite $G$-energy.  The real Banach space
$\Y_G(A,B)$ consists of Hermitian-preserving maps on this domain whose
amplifications have the required continuity on every subnormalized
energy-bounded positive set.  Its norm is
\[
 \|\Phi\|_{\diamond,E}^{G}
 =\sup\bigl\{\|(\Phi\otimes\id_R)(\omega)\|_1:
 \omega\in\mathfrak S(\cH_A\otimes\cH_R),\
 \Tr G\omega_A\leq E\bigr\}.
\]
For discrete $G$, Theorem 3B of the source identifies $\Y_G(A,B)$ with the
ECD completion of the ordinary Hermitian completely bounded maps
$\Y(A,B)$.  Its Theorem 1B identifies $\F_G^0(A,B)$ with the ECD completion
of the cone $\F(A,B)$ of ordinary bounded CP maps.

Here ``discrete'' has the source's meaning: $G$ has finite-multiplicity
eigenvalues tending to $+\infty$.  A zero ground-state energy is allowed.

\section{Proof intuition}

The ECD norm looks weaker than the diamond norm only because high-energy
directions are discounted.  The operator
\[
 R=(I+G/E_0)^{-1/2}
\]
builds that discount into the input.  Passing from $\Phi$ to
$\widehat\Phi(\rho)=\Phi(R\rho R)$ turns the ECD problem into an ordinary
diamond-norm problem, with no loss beyond a factor $2$.  Discreteness of
$G$ is used exactly once: it makes $R$ compact, so the dual of every weighted
map can be approximated by finite spectral corners.

Finite-corner Hermitian cb maps admit normal Wittstock decompositions into
differences of CP maps.  Decomposing a naively chosen approximating sequence
would not suffice, because the positive and negative parts could diverge.
The key bookkeeping step is instead to choose a summably accurate sequence
of corners and decompose its telescoping increments.  The CP norms are then
summable.  On a finite spectral corner $R$ has a bounded inverse, so each CP
piece can be lifted back to an ordinary bounded CP map before taking its ECD
limit.

\section{The weighted norm equivalence}

Fix $E_0>0$ and set
\[
 S=I+G/E_0,\qquad R=S^{-1/2}.
\]
Thus $0<R\leq I$, $\Ran R=\mathcal D(\sqrt G)$, and
$GR^2\leq E_0I$.

\begin{lemma}[Weighted transform]\label{lem:weight}
For every $\Phi\in\Y_G(A,B)$, the formula
\[
 \widehat\Phi(\rho)=\Phi(R\rho R),\qquad \rho\in\T(\cH_A),
\]
defines a Hermitian completely bounded map
$\widehat\Phi:\T(\cH_A)\to\T(\cH_B)$, and
\begin{equation}\label{eq:norm-equivalence}
 \|\widehat\Phi\|_\diamond
 \leq \|\Phi\|_{\diamond,E_0}^{G}
 \leq 2\|\widehat\Phi\|_\diamond .
\end{equation}
The construction commutes with amplification and is injective on
$\Y_G(A,B)$.
\end{lemma}

\begin{proof}
Let $\omega$ be a state on $A$ together with an arbitrary ancilla, and put
\[
 \sigma=(R\otimes I)\omega(R\otimes I).
\]
Then $\sigma\geq0$, $\Tr\sigma\leq1$, and
\[
 \Tr G\sigma_A=\Tr(GR^2)\omega_A\leq E_0.
\]
The subnormalized-state estimate in the proof of Theorem 3A of the source
therefore gives
\[
 \|(\widehat\Phi\otimes\id)(\omega)\|_1
 \leq\|\Phi\|_{\diamond,E_0}^{G}.
\]
For Hermitian maps the diamond norm is the supremum over states.  This proves
that $\widehat\Phi$ extends to a completely bounded map and proves the first
inequality in \eqref{eq:norm-equivalence}.

Conversely, let $\omega$ be a state with $\Tr G\omega_A\leq E_0$.  The
positive quadratic-form expression
\[
 \tau=(S^{1/2}\otimes I)\omega(S^{1/2}\otimes I)
\]
is trace class and satisfies
\[
 \Tr\tau=1+E_0^{-1}\Tr G\omega_A\leq2,
 \qquad (R\otimes I)\tau(R\otimes I)=\omega.
\]
Consequently
\[
 \|(\Phi\otimes\id)(\omega)\|_1
 =\|(\widehat\Phi\otimes\id)(\tau)\|_1
 \leq 2\|\widehat\Phi\|_\diamond.
\]
Taking the supremum proves the second inequality.

For injectivity, suppose $\widehat\Delta=0$.  If
$x\in\mathcal D(\sqrt G)=\Ran R$, then
$|x\rangle\langle x|=R|a\rangle\langle a|R$ for some $a$, so
$\Delta(|x\rangle\langle x|)=0$.  Every finite-energy positive trace-class
operator has a rank-one decomposition whose vectors lie in
$\mathcal D(\sqrt G)$ and whose finite partial sums remain in a common
subnormalized energy-bounded set.  Those sums converge in trace norm, and
the defining continuity of $\Delta\in\Y_G$ gives $\Delta(\rho)=0$.
Linearity on the span $\T_G(\cH_A)$ proves $\Delta=0$.
\end{proof}

\section{Full solution}

\begin{theorem}\label{thm:main}
Let $G$ be a positive discrete unbounded operator on $\cH_A$.  For every
$\Phi\in\Y_G(A,B)$ there are $\Psi_+,\Psi_-\in\F_G^0(A,B)$ such that
\[
 \Phi=\Psi_+-\Psi_-.
\]
Thus the conjecture on page 23 of arXiv:1810.10922 is true.
\end{theorem}

\begin{proof}
Let $T=\widehat\Phi^{,*}:\B(\cH_B)\to\B(\cH_A)$.  This is a normal
Hermitian completely bounded map and
$\|T\|_{\cb}=\|\widehat\Phi\|_\diamond$.

By Theorem 3B of the source, choose
$\Phi_m\in\Y(A,B)$ converging to $\Phi$ in ECD norm.  Lemma~\ref{lem:weight}
shows that $\widehat\Phi_m\to\widehat\Phi$ in diamond norm.  On duals,
\[
 T_m(X)=\widehat\Phi_m^{,*}(X)=R\Phi_m^*(X)R
 \longrightarrow T(X)
 \quad\hbox{in cb norm}.                                      \tag{3}
\]

Let $P_N$ be the spectral projection onto the first $N$ eigenvectors of
$G$, counted with multiplicity, and put $R_N=P_NR=RP_N$.  Since $G$ is
discrete, $R$ is compact and $\|R-R_N\|\to0$.  For fixed $m$,
\[
 T_{m,N}(X)=P_NT_m(X)P_N=R_N\Phi_m^*(X)R_N
\]
is a normal Hermitian cb map with values in the finite corner
$P_N\B(\cH_A)P_N$, and
\[
 \|T_m-T_{m,N}\|_{\cb}
 \leq 2\|R\|\,\|\Phi_m\|_\diamond\,\|R-R_N\|\longrightarrow0. \tag{4}
\]
Equations (3)--(4) show that $T$ is in the cb-norm closure of the normal
Hermitian finite-corner maps.

Choose such maps $S_k$, supported on an increasing sequence of spectral
corners, so that
\[
 \|T-S_k\|_{\cb}\leq4^{-k}.
\]
Set $S_0=0$ and $D_k=S_k-S_{k-1}$.  Then
\begin{equation}\label{eq:summable-D}
 T=\sum_{k=1}^{\infty}D_k
 \quad\hbox{in cb norm},\qquad
 \sum_{k=1}^{\infty}\|D_k\|_{\cb}<\infty.                 
\end{equation}
Each $D_k$ is normal, Hermitian, and supported in some finite spectral
corner $P_{N_k}\B(\cH_A)P_{N_k}$.

The normal Wittstock decomposition theorem~\cite{Paulsen,Wittstock} gives
normal CP maps $C_{k,+},C_{k,-}$ into that same corner such that
\[
 D_k=C_{k,+}-C_{k,-},
 \qquad
 \|C_{k,\pm}\|_{\cb}\leq\|D_k\|_{\cb}+2^{-k}.            \tag{5}
\]
For completeness, the norm control follows from a normal factorization
$D_k(X)=\tfrac12[U^*\pi(X)V+V^*\pi(X)U]$, after balancing $U,V$, and the
identity
\[
 D_k(X)=\tfrac14(U+V)^*\pi(X)(U+V)
       -\tfrac14(U-V)^*\pi(X)(U-V).
\]
An arbitrarily small factorization error yields the $2^{-k}$ allowance.
By \eqref{eq:summable-D}, both series of CP cb norms are summable.

On $P_{N_k}\cH_A$, the inverse
$R_{N_k}^{-1}=P_{N_k}R^{-1}P_{N_k}$ is bounded.  Define a normal CP map
\[
 L_{k,\pm}(X)=R_{N_k}^{-1}C_{k,\pm}(X)R_{N_k}^{-1},
 \qquad X\in\B(\cH_B),
\]
and let $\psi_{k,\pm}:\T(\cH_A)\to\T(\cH_B)$ be its predual.  Thus
$\psi_{k,\pm}$ is an ordinary bounded CP map.  Its weighted dual is exactly
the desired corner piece:
\[
 \widehat\psi_{k,\pm}^{,*}(X)
 =RL_{k,\pm}(X)R=C_{k,\pm}(X).                              \tag{6}
\]
Lemma~\ref{lem:weight}, (5), and (6) give
\[
 \sum_k\|\psi_{k,\pm}\|_{\diamond,E_0}^{G}
 \leq2\sum_k\|C_{k,\pm}\|_{\cb}<\infty.                 \tag{7}
\]

The partial sums $\sum_{k\leq n}\psi_{k,\pm}$ are ordinary bounded CP
maps.  By (7) they converge in ECD norm, and Theorem 1B of the source says
that their limits $\Psi_\pm$ belong to $\F_G^0(A,B)$.  Taking weighted
duals and using \eqref{eq:summable-D} and (6) yields
\[
 \widehat\Phi^{,*}=T
 =\sum_k(C_{k,+}-C_{k,-})
 =\widehat\Psi_+^{,*}-\widehat\Psi_-^{,*}.
\]
Hence $\widehat\Phi=\widehat\Psi_+-\widehat\Psi_-$.  Injectivity in
Lemma~\ref{lem:weight} finally gives
$\Phi=\Psi_+-\Psi_-$ on $\T_G(\cH_A)$.
\end{proof}

\section{Checks, scope, and novelty guardrail}

The factor $2$ in Lemma~\ref{lem:weight} comes exactly from
$\Tr[(I+G/E_0)\omega_A]\leq2$; no finite-dimensional approximation is
hidden in that step.  The only use of discreteness is compactness of $R$ in
(4).  The proof does not claim the conjecture for a positive $G$ with
noncompact resolvent.

The telescoping construction is essential: it turns convergence of
Hermitian corners into absolute summability of their CP parts.  Normality is
also essential because it ensures that every lifted CP map has a predual on
trace class.  The finite-corner inverse is always bounded, even though
$R^{-1}$ is unbounded on the full Hilbert space.

On 11 August 2026, bounded novelty searches covered the run's registry,
solution, attempt, and proof-gap indexes; the exact source identifier and
notation $\Y_G(A,B)$ and $\F_G^0(A,B)$; and arXiv/web-facing variants of
``ECD completion difference of completely positive maps'', ``Hermitian
energy-constrained diamond norm Wittstock'', and the exact conjecture text.
They found the source, general completely bounded/decomposable-map
literature, and an already-recorded answer to the source's separate operator
$E$-norm question, but no later resolution of this conjecture or this
weighted compact-corner argument.  This is a bounded novelty guardrail, not
a claim of priority.

\paragraph{Human-review recommendation.}
Review as a candidate full solution.  The three highest-value checks are:
(i) the quadratic-form identity used for $\tau$ in
Lemma~\ref{lem:weight}; (ii) the normal, norm-controlled Wittstock
decomposition in finite corners; and (iii) the passage from vanishing on
weighted rank-one operators to all of $\T_G(\cH_A)$ in the injectivity
argument.

\begin{thebibliography}{9}
\bibitem{Shirokov}
M.~E.~Shirokov,
\emph{On completion of the cone of CP linear maps with respect to the
energy-constrained diamond norm},
J. Math. Phys. \textbf{60} (2019), 102202; arXiv:1810.10922v2.

\bibitem{Paulsen}
V.~Paulsen,
\emph{Completely Bounded Maps and Operator Algebras},
Cambridge Studies in Advanced Mathematics 78, Cambridge University Press,
2002.

\bibitem{Wittstock}
G.~Wittstock,
\emph{Ein operatorwertiger Hahn--Banach Satz},
J. Funct. Anal. \textbf{40} (1981), 127--150.
\end{thebibliography}

\end{document}
